\documentclass[10pt,a4paper]{amsart}
\usepackage[margin=19mm]{geometry}
\usepackage{amsmath,amssymb,mathtools}
\usepackage[scr=rsfs,cal=euler]{mathalfa}
\usepackage{xfrac}
\usepackage[unicode,hidelinks,bookmarks=false]{hyperref}
\newcommand{\ie}{i.e.\ }
\newcommand{\cf}{cf.\ }
\newcommand{\resp}{respectively\ }
\newcommand{\Subsection}{\subsection}
\providecommand{\nobreakdash}{\nobreak\hbox{-}}
\setlength{\emergencystretch}{2em}
\multlinegap0pt
\usepackage{bm}
\usepackage{bbm}
\usepackage{dsfont}
\usepackage{xspace}
\usepackage{cancel}
\usepackage{mathtools}
\usepackage{amsthm}
\makeatletter
\@ifundefined{theorem}{\newtheorem{theorem}{Theorem}}{}
\@ifundefined{lemma}{\newtheorem{lemma}[theorem]{Lemma}}{}
\makeatother
\makeatletter
\newcommand{\R}{\mathbb{R}}
\newcommand{\conv}[1]{\mathsf{conv}\left(#1\right)}
\newcommand{\defi}{\stackrel{\mathrm{\scriptscriptstyle def}}{=}}
\newcommand{\faces}[1]{{\mathsf{faces}^\ast}\left(#1\right)}
\newcommand{\mP}{\mathcal{P}}
\expandafter\ifx\csname proof\endcsname\relax
\renewenvironment{proof}%
{%
 \par\noindent{\bfseries\upshape \proofname\ }%
}%
{\jmlrQED}
\fi
\newcommand{\verts}[1]{\mathcal{V}_{#1}}
\makeatother
\title[Linear Convergence of the Frank-Wolfe Algorithm over Product]{Linear Convergence of the Frank-Wolfe Algorithm over Product Polytopes}
\author{Gabriele Iommazzo, David Mart\'inez-Rubio, Francisco Criado, Elias Wirth, Sebastian Pokutta}
\date{Source: arXiv:2505.11259v2 (CC BY 4.0)}
\begin{document}
\maketitle

\noindent\emph{Source and licence.} Linear Convergence of the Frank-Wolfe Algorithm over Product Polytopes. Version arXiv:2505.11259v2 (CC BY 4.0), distributed under the Creative Commons Attribution 4.0 International licence, as recorded in the arXiv licence metadata for this exact version and shown on the abstract page https://arxiv.org/abs/2505.11259v2. Full rights notice: licenses/arxiv-2505.11259-CC-BY-4.0.txt.\par\smallskip

\noindent\textit{Editorial scope.} Counted unit: the Lemma whose source label is \texttt{fact:cart\_prod\_kpolytopes} in the article (source0.tex lines 1192--1196, source label \texttt{fact:cart\_prod\_kpolytopes}), delivered together with the article's own complete original proof of it (source0.tex lines 1198--1244, one \texttt{proof} environment of 47 lines). No mathematics is written, repaired or re-derived: statement and proof are the article's own text line for line. Required context retained uncounted: none. Every definition, display and label of those lines is retained unchanged. Omitted from this package and not redistributed: 1--1191, 1197--1197, 1245--2292. Nothing inside a retained excerpt is omitted or altered: every retained body is delivered exactly as the article's lines are written, so no omission receipt of any kind accompanies this package. Citations: the retained excerpts carry no citation command at all, so no bibliography entry of the article is transcribed and no reference is redistributed. Reference adaptation: none was needed and none was made; every reference inside the delivered unit resolves inside it, so no unresolved reference or citation is produced. Printed slips kept literal: no printed word, symbol, hypothesis or number of the counted statement or of its proof was altered. Layout and class: the article's own class is \texttt{article} with packages including neurips\_2025, inputenc, fontenc, url, microtype, textcomp, lmodern, babel, lipsum, graphicx, adjustbox, scalerel, stackengine, amsmath; the wrapper is the lane's pinned \texttt{amsart} editorial wrapper at 10pt on A4, which restates the theorem-like environments the retained text needs and adds the article's own macro definitions that the retained text needs (\texttt{\textbackslash R}, \texttt{\textbackslash conv}, \texttt{\textbackslash defi}, \texttt{\textbackslash faces}, \texttt{\textbackslash mP}, \texttt{\textbackslash verts}), with extra wrapper packages (bm, bbm, dsfont, xspace, cancel, mathtools). The delivered numbering is the unit's own. Assets: no retained span contains a figure environment, a graphics-inclusion command, a PDF-page inclusion, a table or a TikZ picture; no image file is copied into this package, the package declares no asset dependency and no image file is redistributed with it. Licence: version arXiv:2505.11259v2 of this article is distributed under the Creative Commons Attribution 4.0 International licence, as recorded in the arXiv licence metadata for this exact version and shown on the abstract page https://arxiv.org/abs/2505.11259v2. A full rights notice is delivered at licenses/arxiv-2505.11259-CC-BY-4.0.txt. Characters: every retained excerpt is pure ASCII, so the delivered engine, XeLaTeX with its default Latin Modern text font, reports no missing character.
\begin{lemma}[\textbf{Cartesian product of $k$ polytopes}]\label{fact:cart_prod_kpolytopes}
The Cartesian product $\mP = \Pi_{i \in [k]} \mP_i$ of $k$ polytopes, each in $\R^n$ for some $n \in \mathbb{N}_{>0}$, is a polytope of dimension $kn$.
In particular, $\verts{\mP} = \Pi_{i \in [k]}\verts{\mP_i}$.
Further, if $F_\mP \in \faces{\mP}$, then  $F_\mP$ is the product of faces of all the individual polytopes, i.e., $\forall i \in [k]$ there exist $F_{\mP_i} \in \faces{\mP_i}$ such that $F_\mP = \Pi_{i \in [k]} F_{\mP_i}$.
\end{lemma}

\begin{proof}
For all $i \in [k]$, we denote the $i$-th polytope by $\mP_i \defi \{x \in \R^n : A_{\mP_i} x \le b_{\mP_i}\}$, where $A_{\mP_i} \in \R^{m_{\mP_i} \times n}$ and $b_{\mP_i} \in \R^{m_{\mP_i}}$.
Compactness and convexity are inherently preserved in the Cartesian product, so
\[
\begin{array}{rcl}
\Pi_{i \in [k]} &=& \{(x^1, \dots, x^k) \in \R^n\times\cdots \times\R^n\,:\, x^1 \in \mP_1, \dots, x^k \in \mP_k\}\\
&=& \{(x \in \R^{kn}\,:\, \forall i \in [k]\,,\, A_{\mP_i} x^i \le b_{\mP_i}\}\,.
\end{array}
\]
By linearity of the constraints and independence of each block $i \in [k]$,
%of constraints only involves one polytope at a time,
% \[
% \mP_1 \times \cdots \times \mP_k = \left\{x \in \R^{kn} \,:\:
% \left(\begin{array}{cccc}
% A_{\mP_1} & 0 & \cdots & 0 \\
% 0 & A_{\mP_2} & & \vdots \\
% \vdots & & \ddots & 0 \\
% 0 & \cdots & 0 & A_{\mP_k} \\
% \end{array}\right)
% \begin{pmatrix}
%     x_1\\
%     \vdots\\
%     x_k
% \end{pmatrix} \le \begin{pmatrix}
%     b_{\mP_1} \\
%     \vdots\\
%     b_{\mP_k}
% \end{pmatrix}\right\}\,,
% \]
the product is still a polytope described by a constraint matrix in $\R^{(m_{\mP_1}+\dots+m_{\mP_k})\times kn}$ with diagonal blocks $A_{\mP_i}$, for $i \in [k]$, and right-hand side coefficients $b_{\mP_i} \in \R^{m_{\mP_1}+\dots+m_{\mP_k}}$.
Moreover, since $\forall i \in [k]$, $\mP_i = \conv{\mP_i}$ and $\conv{\Pi_{i \in [k]}\mP_i} = \Pi_{i \in [k]}\conv{\mP_i}$, then $\verts{\mP} = \Pi_{i \in [k]}\verts{\mP_i}$.
% We show that the lemma is valid for two polytopes $\mP$ and $\mQ$.
% Denote by $\text{conv}(\verts(\mP)$ and $\text{conv}(\verts(\mQ)$ the convex hull of $\mP$ and $\mQ$, respectively, and by $\Delta_{n-1}$, is given by $\Delta_{n-1} = \{z \in \R^n : z_i \ge 0, \text{ for } i \in [n],\, \|z\|_1 = 1\}$ the standard $(n-1)$-simplex.
% We use the definition of polytopes as the convex hull of their vertices, i.e., $\mP = \text{conv}(\verts(\mP)$ and $\mQ = \text{conv}(\verts(\mQ)$, and $\mP \times \mQ = \text{conv}(\verts(\mP) \times \text{conv}(\verts(\mQ)$. So
% \begin{equation}
% (x,y) = \left(\sum_{i=1}^{|\verts(\mP|} \lambda_i v_i, \sum_{j=1}^{|\verts(\mQ|} \mu_j w_j\right) \in \text{conv}(\verts(\mP) \times \text{conv}(\verts(\mQ)\,,
% \label{eq:convconv_mPmQ}
% \end{equation}
% for $v_i \in \verts(\mP$, $w_j \in \verts(\mQ$, $\lambda, \mu \in \Delta_{n-1}$. We will show that $\text{conv}(\verts(\mP) \times \text{conv}(\verts(\mQ) = \text{conv}\left(\verts(\mP \times \verts(\mQ\right)$.
% If $(x',y') \in \text{conv}\left(\verts(\mP \times \verts(\mQ\right)$, then $(x',y') = \sum_{i=1}^{|\verts(\mP|}\sum_{j=1}^{|\verts(\mQ|} \gamma_{ij} (v_i, w_j)$, where $\gamma \in \Delta_{|\verts(\mP| + |\verts(\mQ| - 1}$.
% Expanding and summing over the appropriate indices yields $(x',y') = (\sum_{i=1}^{|\verts(\mP|}v_i \sum_{j=1}^{|\verts(\mQ|}\gamma_{ij}, \sum_{j=1}^{|\verts(\mQ|}w_j \sum_{i=1}^{|\verts(\mP|}\gamma_{ij})$, where each $\sum_{i=1}^{|\verts(\mP|}\sum_{j=1}^{|\verts(\mQ|}\gamma_{ij}$ sums up to 1, thus contributing to a convex combination within $\mP$ and $\mQ$, based on their respective indices. Renaming the sums as $\sum_{j=1}^{|\verts(\mQ|}\gamma_{ij} = \lambda_i$ and $\sum_{i=1}^{|\verts(\mP|}\gamma_{ij} = \mu_j$ for clarity, we get \cref{eq:convconv_mPmQ} indicating that $\text{conv}(\verts(\mP \times \verts(\mQ)$ is equivalent to $\text{conv}(\verts(\mP) \times \text{conv}(\verts(\mQ)$, thus completing the proof.
% Extending this proof to encompass $k$ polytopes employs the same steps, albeit with heavier notation, so we skip it here for ease of readability.
Finally, a polytope face is the set of points optimizing some linear functional over the polytope.
Consider the linear functional defined by $n = (n^1, \ldots, n^k) \in \R^k$.
A point $(p^1, \ldots, p^k) \in \mP$ optimizes the linear functional defined by $n$ if and only if, each $p^i$ optimizes the linear functional defined by the corresponding $n^i$ over $\mP_i$, as a linear program in the product is solved in each component independently.
If, for all $i \in [k]$, we denote the faces defined by that condition as $F_i \subseteq \mP_i$, the set of points optimizing $n$,is exactly $F_1 \times F_2 \times \cdots \times F_k \subset \mP_1 \times \mP_2 \times \cdots \times \mP_k$, i.e., a face $F_\mP$ of $\mP$.
\end{proof}

\end{document}
